\documentclass[11pt,a4paper]{article}

\usepackage[utf8]{inputenc}
\usepackage[T1]{fontenc}
\usepackage[brazil]{babel}
\usepackage{lmodern}
\renewcommand{\familydefault}{\sfdefault}
\usepackage{microtype}
\usepackage{needspace}
\usepackage{etoolbox}
\usepackage[a4paper, margin=2.4cm, headheight=14pt]{geometry}
\usepackage{xcolor}
\usepackage{listings}
\usepackage{tabularx}
\usepackage{array}
\usepackage{booktabs}
\usepackage{enumitem}
\usepackage{amsmath, amssymb}
\usepackage[explicit]{titlesec}
\usepackage{fancyhdr}
\usepackage{lastpage}
\usepackage[most]{tcolorbox}
\usepackage{tikz}
\usetikzlibrary{arrows.meta, positioning, shapes.geometric, calc, fit, backgrounds}
\usepackage[hidelinks]{hyperref}

% ------------------------------------------------------------
% Cores do tema metropolis, as mesmas dos slides
% ------------------------------------------------------------
\definecolor{mDarkTeal}{HTML}{23373B}
\definecolor{mLightBrown}{HTML}{EB811B}
\definecolor{mLightGreen}{HTML}{14B03D}
\definecolor{mBackground}{HTML}{FAFAFA}
\definecolor{termbg}{RGB}{40,44,52}
\definecolor{termfg}{RGB}{220,220,220}
\definecolor{codekw}{RGB}{235,129,27}
\definecolor{codecom}{RGB}{140,150,160}
\definecolor{codestr}{RGB}{152,195,121}

\hypersetup{pdftitle={Processos e comunicação entre processos}, pdfauthor={Weslei Ferreira Santos}}

\setlength{\parindent}{0pt}
\setlength{\parskip}{0.55em}
\setlist{itemsep=0.2em, topsep=0.3em}
\setlist[itemize]{label=\textcolor{mDarkTeal}{\small$\blacktriangleright$}}

% ------------------------------------------------------------
% Títulos no estilo dos slides
% ------------------------------------------------------------
\newcommand{\barradesecao}[1]{%
  \begin{tcolorbox}[colback=mDarkTeal, colframe=mDarkTeal, sharp corners, boxrule=0pt,
      left=10pt, right=10pt, top=7pt, bottom=7pt, coltext=white, fontupper=\Large\bfseries]
  \thesection\quad #1
  \end{tcolorbox}}
\titleformat{\section}[block]{}{}{0pt}{\barradesecao{#1}}
\titleformat{name=\section, numberless}[block]{}{}{0pt}{%
  \begin{tcolorbox}[colback=mDarkTeal, colframe=mDarkTeal, sharp corners, boxrule=0pt,
      left=10pt, right=10pt, top=7pt, bottom=7pt, coltext=white, fontupper=\Large\bfseries]
  #1
  \end{tcolorbox}}
\titlespacing*{\section}{0pt}{1.2em}{0.6em}

\titleformat{\subsection}[block]{\large\bfseries\color{mDarkTeal}}{}{0pt}%
  {\thesubsection\quad #1}[\vspace{-\parskip}\vspace{3pt}{\color{mLightBrown}\rule{3.2cm}{1.4pt}}]
\titleformat{name=\subsection, numberless}[block]{\large\bfseries\color{mDarkTeal}}{}{0pt}%
  {#1}[\vspace{-\parskip}\vspace{3pt}{\color{mLightBrown}\rule{3.2cm}{1.4pt}}]
\titlespacing*{\subsection}{0pt}{1.1em}{0.4em}
\pretocmd{\subsection}{\needspace{6\baselineskip}}{}{\PackageWarning{apostila}{pretocmd falhou}}
\pretocmd{\section}{\needspace{10\baselineskip}}{}{\PackageWarning{apostila}{pretocmd falhou}}

\titleformat{\subsubsection}[runin]{\bfseries\color{mDarkTeal}}{}{0pt}{#1.}
\titlespacing*{\subsubsection}{0pt}{0.8em}{0.6em}

% ------------------------------------------------------------
% Cabeçalho e rodapé
% ------------------------------------------------------------
\pagestyle{fancy}
\fancyhf{}
\fancyhead[L]{\small\color{mDarkTeal} MATA58 · Sistemas Operacionais}
\fancyhead[R]{\small\color{mDarkTeal} Processos e comunicação entre processos}
\fancyfoot[R]{\small\color{mDarkTeal}\thepage/\pageref{LastPage}}
\renewcommand{\headrule}{\color{mLightBrown}\hrule height 0.8pt}
\renewcommand{\footrulewidth}{0pt}

% ------------------------------------------------------------
% Blocos no estilo "block=fill" do metropolis
% ------------------------------------------------------------
\tcbset{
  bloco/.style={enhanced, breakable, sharp corners, boxrule=0pt, frame hidden,
    colback=black!4, colbacktitle=black!11, fonttitle=\bfseries,
    left=8pt, right=8pt, top=5pt, bottom=6pt, toptitle=3pt, bottomtitle=3pt,
    before skip=0.9em, after skip=0.9em}
}
\newtcolorbox{definicao}[1][Definição]{bloco, coltitle=mDarkTeal, title={#1}}
\newtcolorbox{alerta}[1][Atenção]{bloco, coltitle=mLightBrown, title={#1}}
\newtcolorbox{exemplo}[1][Exemplo]{bloco, coltitle=mLightGreen!80!black, title={#1}}
\newtcolorbox{pergunta}[1][Pare e pense]{bloco, colback=mLightBrown!7, colbacktitle=mLightBrown!20,
  coltitle=mLightBrown!70!black, title={#1}}

% ------------------------------------------------------------
% Código e terminal, iguais aos slides
% ------------------------------------------------------------
\lstdefinestyle{terminal}{
    backgroundcolor=\color{termbg},
    basicstyle=\ttfamily\small\color{termfg},
    breaklines=true,
    frame=single,
    rulecolor=\color{termbg},
    framerule=3pt,
    xleftmargin=4pt,
    xrightmargin=4pt,
    columns=fullflexible,
    keepspaces=true,
    showstringspaces=false,
    aboveskip=0.8em,
    belowskip=0.8em
}
\lstdefinestyle{codigo}{
    style=terminal,
    language=C,
    basicstyle=\ttfamily\footnotesize\color{termfg},
    keywordstyle=\color{codekw}\bfseries,
    commentstyle=\color{codecom}\itshape,
    stringstyle=\color{codestr},
    morekeywords={pid_t, size_t, ssize_t},
    moredelim=[is][\color{codekw}\bfseries]{@}{@}
}
\lstdefinestyle{arquivo}{
    style=codigo,
    numbers=left,
    numberstyle=\tiny\color{codecom},
    numbersep=8pt,
    xleftmargin=22pt,
    framexleftmargin=18pt
}

% ------------------------------------------------------------
% Figuras
% ------------------------------------------------------------
\tikzset{
    proc/.style={draw=mDarkTeal, fill=mDarkTeal!8, rounded corners=3pt, thick,
        minimum width=2.1cm, minimum height=0.8cm, align=center, font=\small},
    procf/.style={proc, draw=mDarkTeal!70, fill=white, minimum width=1.2cm, font=\footnotesize},
    destaque/.style={proc, draw=mLightBrown, fill=mLightBrown!12},
    arquivo/.style={draw=black!50, fill=black!5, thick, minimum width=2.2cm,
        minimum height=1.1cm, align=center, font=\small},
    seta/.style={-{Stealth[length=2.2mm]}, thick, draw=mDarkTeal},
    setad/.style={-{Stealth[length=2.2mm]}, very thick, draw=mLightBrown},
    rotulo/.style={font=\scriptsize, fill=white, inner sep=1.5pt},
    celula/.style={draw=mDarkTeal, minimum size=0.7cm, font=\ttfamily\small, fill=white},
    regiao/.style={draw=mDarkTeal, fill=mDarkTeal!6, minimum width=3.6cm, minimum height=0.75cm,
        font=\small, align=center}
}

\newcommand{\raiz}{../../../..}
\newcommand{\cmd}[1]{\texttt{#1}}

\begin{document}

% ============================================================
% Capa no estilo da página de título dos slides
% ============================================================
\begin{titlepage}
\thispagestyle{empty}
\vspace*{5cm}
{\color{mDarkTeal}\fontsize{28}{34}\selectfont\bfseries Processos e comunicação\\[0.2em] entre processos\par}
\vspace{0.8em}
{\Large\color{mDarkTeal!80} fork, memória independente e pipe na prática\par}
\vspace{1em}
{\color{mLightBrown}\rule{\linewidth}{1.2pt}}
\vspace{1em}

{\large Apostila da aula de 5 de outubro de 2026\par}
\vspace{3em}
{\large Weslei Ferreira Santos\par}
\vspace{0.6em}
MATA58 · Sistemas Operacionais\\
Instituto de Computação / UFBA

\vfill
{\small\color{black!60} Material-base: Prof. Maycon Leone M. Peixoto. Os programas citados estão no repositório da disciplina e podem ser compilados com \cmd{make}.\par}
\end{titlepage}

\tableofcontents
\clearpage

% ============================================================
\section{Antes de começar}
% ============================================================

Esta apostila acompanha a aula sobre processos e comunicação entre processos. Os slides mostram o caminho em poucas palavras; aqui cada ideia aparece com mais calma, com os programas completos, as saídas esperadas e os detalhes que costumam gerar dúvida.

A aula inteira gira em torno de uma pergunta:

\begin{center}
{\large\bfseries\color{mDarkTeal} Como um processo recebe o resultado calculado por outro?}
\end{center}

Para respondê-la, seguimos uma ordem deliberada. Primeiro entendemos o que é um processo e como ele nasce com \cmd{fork()}. Depois constatamos, com um experimento, que pai e filho \emph{não} compartilham a memória. Só então faz sentido apresentar o \cmd{pipe}, o mecanismo que o sistema operacional oferece para transportar dados de um processo a outro.

\subsection{O que você deve conseguir ao final}

\begin{itemize}
    \item Identificar os valores que \cmd{fork()} devolve ao pai e ao filho e explicar de onde cada um continua a execução.
    \item Explicar por que uma variável alterada no filho não muda no pai.
    \item Desenhar a árvore de processos criada por um laço que chama \cmd{fork()}.
    \item Diferenciar \emph{esperar o término} de um filho (\cmd{waitpid}) de \emph{receber dados} dele (\cmd{pipe}).
    \item Completar um programa em que o filho envia um inteiro ao pai por um pipe.
\end{itemize}

\subsection{Como usar os exemplos}

Todos os programas citados ficam no repositório da disciplina. Em um terminal Linux, na raiz do repositório:

\begin{lstlisting}[style=terminal]
$ make test            # compila os exemplos e confere o comportamento
$ make run-fork-basico # retornos de fork
$ make run-memoria     # mesmo endereco, valores diferentes
$ make run-fork        # fork2.c
$ make run-pipe-soma   # filho envia 30 ao pai
$ make run-pipe        # pipe2.c: raizes da equacao de 2o grau
\end{lstlisting}

Cada \cmd{make run-...} primeiro imprime o comando que vai executar (por exemplo, \cmd{./build/fork2}) e só depois a saída do programa. Na primeira execução aparece antes também a linha do \cmd{gcc} que compila o exemplo. As saídas mostradas nesta apostila incluem a linha do comando, mas não a do \cmd{gcc}.

Os exemplos são compilados com \cmd{-std=c11 -Wall -Wextra -Wpedantic -Werror}: qualquer aviso do compilador é tratado como erro. Rode cada programa enquanto lê a seção correspondente. Os números de PID das saídas mostradas aqui vão ser diferentes na sua máquina, e isso é esperado.

% ============================================================
\section{Programa e processo}
% ============================================================

\subsection{Duas palavras que parecem sinônimas}

Um \textbf{programa} é um arquivo: uma sequência de instruções e dados guardada no disco. Ele não faz nada sozinho. Quando pedimos ao sistema operacional para executá-lo, o SO cria um \textbf{processo}.

\begin{definicao}
Um \textbf{processo} é um programa em execução, junto com todo o estado necessário para executá-lo: o identificador (PID), o conteúdo dos registradores, a memória (código, dados, heap e pilha), os arquivos abertos e outras informações mantidas pelo núcleo do SO.
\end{definicao}

A diferença fica clara quando executamos o mesmo programa duas vezes, em dois terminais. Há um único arquivo e dois processos. Cada um tem o seu PID, a sua pilha, o seu heap e os seus registradores.

\begin{center}
\begin{tikzpicture}
    \node[arquivo] (prog) {\texttt{fork2}\\\scriptsize arquivo no disco};
    \node[proc, right=3cm of prog, yshift=0.75cm] (p1) {Processo\\\scriptsize PID 4210};
    \node[proc, right=3cm of prog, yshift=-0.75cm] (p2) {Processo\\\scriptsize PID 4231};
    \draw[seta] (prog.east) -- node[rotulo, above, sloped] {executa} (p1.west);
    \draw[seta] (prog.east) -- node[rotulo, below, sloped] {executa} (p2.west);
    \node[right=0.4cm of p1, font=\scriptsize, align=left] {pilha, heap,\\registradores};
    \node[right=0.4cm of p2, font=\scriptsize, align=left] {pilha, heap,\\registradores};
\end{tikzpicture}

{\small Os PIDs da figura são ilustrativos.}
\end{center}

\subsection{O que o SO guarda sobre cada processo}

Para cada processo, o núcleo mantém um registro (frequentemente chamado de \emph{bloco de controle do processo}, ou PCB). Entre as informações guardadas estão:

\begin{itemize}
    \item \textbf{PID}: o número que identifica o processo enquanto ele existe.
    \item \textbf{PPID}: o PID do processo que o criou (o pai).
    \item \textbf{Estado}: executando, pronto, bloqueado esperando algo, terminado.
    \item \textbf{Contexto}: o valor dos registradores, incluindo o contador de programa, salvo quando o processo sai da CPU.
    \item \textbf{Espaço de endereçamento}: a memória que o processo enxerga.
    \item \textbf{Tabela de descritores}: os arquivos, pipes e terminais que o processo tem abertos.
\end{itemize}

No Linux, essas informações aparecem em comandos como \cmd{ps -f} e \cmd{pstree -p}, vistos na aula de introdução ao Linux.

\subsection{O espaço de endereçamento}

Cada processo enxerga a sua memória como se fosse a única do computador. Essa memória está organizada em regiões:

\begin{center}
\begin{tikzpicture}[node distance=0pt]
    \node[regiao] (pilha) {pilha};
    \node[regiao, below=of pilha, fill=white, minimum height=1.1cm] (livre) {\scriptsize $\downarrow$ \quad espaço livre \quad $\uparrow$};
    \node[regiao, below=of livre] (heap) {heap};
    \node[regiao, below=of heap] (dados) {dados globais};
    \node[regiao, below=of dados] (texto) {código (texto)};
    \node[right=0.5cm of pilha, font=\footnotesize, align=left] {variáveis locais e chamadas de função};
    \node[right=0.5cm of heap, font=\footnotesize, align=left] {\texttt{malloc}, \texttt{calloc}};
    \node[right=0.5cm of dados, font=\footnotesize, align=left] {variáveis globais e \texttt{static}};
    \node[right=0.5cm of texto, font=\footnotesize, align=left] {instruções do programa};
    \node[left=0.3cm of pilha, font=\scriptsize, color=black!60] {endereços altos};
    \node[left=0.3cm of texto, font=\scriptsize, color=black!60] {endereços baixos};
\end{tikzpicture}
\end{center}

A figura é uma \textbf{visão simplificada}. Um processo real também tem regiões para bibliotecas compartilhadas e para arquivos mapeados em memória (\cmd{mmap}), e o Linux sorteia a posição de várias regiões a cada execução (ASLR, \emph{Address Space Layout Randomization}), o que dificulta ataques que dependem de endereços fixos. Para esta aula basta a organização geral: código, dados, heap e pilha.

Os endereços que o programa usa são \textbf{virtuais}: o hardware e o SO traduzem cada endereço virtual para uma posição física da memória, e essa tradução é diferente para cada processo. Esse detalhe vai explicar, na Seção~\ref{sec:memoria}, por que o mesmo endereço pode guardar valores diferentes em dois processos.

\begin{pergunta}
Se uma execução de um programa altera uma variável, a outra execução, em outro terminal, necessariamente observa a alteração? Anote sua resposta. Vamos testá-la nas próximas seções.
\end{pergunta}

% ============================================================
\section{Criando processos com fork()}
% ============================================================

\subsection{Uma chamada, dois retornos}

No Linux e em outros sistemas Unix, um processo cria outro com a chamada de sistema \cmd{fork()}. Ela duplica o processo que a chamou: o original passa a ser o \textbf{pai} e a cópia é o \textbf{filho}.

\begin{lstlisting}[style=codigo]
#include <unistd.h>

pid_t fork(void);
\end{lstlisting}

O aspecto mais estranho de \cmd{fork()} é que ela é chamada uma vez e \textbf{retorna duas vezes}: uma no pai e outra no filho. O valor de retorno é o que permite a cada processo saber quem ele é.

\begin{center}
\begin{tabular}{@{}lll@{}}
\toprule
\textbf{Onde} & \textbf{Valor devolvido} & \textbf{Significado} \\
\midrule
No filho & \cmd{0} & ``Eu sou o filho.'' \\
No pai & PID do filho (positivo) & ``Criei este filho.'' \\
No pai, em caso de falha & \cmd{-1} & Nenhum filho foi criado; \cmd{errno} diz o motivo. \\
\bottomrule
\end{tabular}
\end{center}

Por isso o código que usa \cmd{fork()} quase sempre tem a mesma forma:

\begin{lstlisting}[style=codigo]
pid_t pid = fork();
if (pid == -1) {
    /* falha: nenhum filho foi criado nesta chamada */
} else if (pid == 0) {
    /* este caminho executa no filho */
} else {
    /* este caminho executa no pai; pid identifica o filho */
}
/* pai e filho podem chegar aqui */
\end{lstlisting}

\begin{alerta}[O filho não recomeça o main]
O filho é uma cópia do pai \emph{no ponto da chamada}. Ele não executa as linhas anteriores a \cmd{fork()}: continua logo depois dela, exatamente como o pai, só que recebendo \cmd{0} como retorno. Variáveis criadas antes da chamada já existem no filho, com os mesmos valores que tinham no pai naquele instante.
\end{alerta}

\subsection{O que o filho herda}

Logo após um \cmd{fork()} bem-sucedido, o filho tem:

\begin{itemize}
    \item uma \textbf{cópia} de toda a memória do pai: código, variáveis globais, heap e pilha;
    \item os mesmos valores nos registradores, inclusive o ponto do programa onde a execução continua;
    \item cópias dos \textbf{descritores de arquivo} abertos, que apontam para os mesmos arquivos e pipes que os do pai (isso será essencial na Seção~\ref{sec:pipe});
    \item um PID \textbf{novo} e, como PPID, o PID do pai.
\end{itemize}

Copiar toda a memória a cada \cmd{fork()} seria caro. O Linux usa a técnica de \textbf{cópia na escrita} (\emph{copy-on-write}): logo após a chamada, pai e filho apontam para as mesmas páginas físicas, marcadas como somente leitura. Quando um dos dois tenta escrever numa página, o SO faz uma cópia dela só para quem escreveu. Para o programador o efeito é o mesmo de uma cópia completa: \textbf{o que um escreve, o outro não vê}.

\subsection{getpid() e getppid()}

Duas chamadas simples ajudam a acompanhar quem é quem:

\begin{itemize}
    \item \cmd{getpid()} devolve o PID do processo que está executando a chamada;
    \item \cmd{getppid()} devolve o PID do pai desse processo.
\end{itemize}

Note a diferença: no pai, o valor devolvido por \cmd{fork()} é o PID do \emph{filho}, e não o do próprio pai.

\subsection{Um exemplo completo}

O programa \cmd{examples/dia05/fork\_basico.c} mostra os dois retornos lado a lado:

\lstinputlisting[style=arquivo]{\raiz/examples/dia05/fork_basico.c}

Uma execução possível:

\begin{lstlisting}[style=terminal]
$ make run-fork-basico
./build/fork-basico
Antes do fork: PID 2020435
Filho: fork retornou 0; meu PID 2020436; meu pai 2020435
Pai: fork retornou 2020436; meu PID 2020435
Pai: filho 2020436 terminou com codigo 0
\end{lstlisting}

Repare que o número devolvido ao pai (\cmd{2020436}) é o PID que o filho obtém com \cmd{getpid()}, e que o \cmd{getppid()} do filho é o PID do pai.

\subsection{Quem executa primeiro?}

Depois de \cmd{fork()}, pai e filho são dois processos independentes, disputando a CPU. O escalonador do SO decide quem roda, em que ordem e por quanto tempo. Em computadores com vários núcleos, os dois podem até rodar ao mesmo tempo.

\begin{alerta}[Não conte com a ordem]
Na saída acima, a linha do filho apareceu antes da linha do pai. Em outra execução a ordem pode se inverter. Um programa correto não pode depender disso: se a ordem importa, ela precisa ser imposta com sincronização, como \cmd{waitpid} ou um pipe.
\end{alerta}

\subsection{Um detalhe de C: buffers de saída}

As funções de \cmd{stdio.h}, como \cmd{printf}, não escrevem diretamente no terminal. Elas acumulam o texto num \emph{buffer} na memória do processo e só o entregam ao SO quando o buffer enche, quando aparece uma quebra de linha (se a saída é um terminal) ou quando o programa termina normalmente.

Esse buffer faz parte da memória e, portanto, é copiado pelo \cmd{fork()}. Duas consequências:

\begin{enumerate}
    \item Texto que ainda estava no buffer antes do \cmd{fork()} pode ser impresso \textbf{duas vezes}, uma por processo. É por isso que \cmd{fork\_basico.c} chama \cmd{fflush(stdout)} antes de \cmd{fork()}.
    \item A função \cmd{\_exit} encerra o processo imediatamente, \textbf{sem} esvaziar esses buffers. Se o filho imprime algo e chama \cmd{\_exit}, precisa chamar \cmd{fflush(stdout)} antes, ou o texto pode se perder. Isso acontece de fato quando a saída é redirecionada para um arquivo ou para outro programa.
\end{enumerate}

Por que então o filho usa \cmd{\_exit} e não \cmd{exit}? Porque \cmd{exit} esvaziaria também os buffers \emph{herdados} do pai e executaria funções registradas pelo pai com \cmd{atexit}. Em um filho, \cmd{\_exit} encerra apenas o que é dele.

Isso não quer dizer que \cmd{return} ou \cmd{exit} no filho sejam sempre um erro. No \cmd{fork2.c} (Seção~4), todos os processos terminam com \cmd{return} ao final de \cmd{main}, e está correto: nenhum texto fica no buffer antes dos \cmd{fork()}, e só o processo original imprime. O cuidado com \cmd{\_exit} importa quando o filho termina \emph{no meio} do código do pai, como nos exemplos com pipe, ou quando há saída pendente no buffer.

% ============================================================
\section{fork() dentro de um laço}
% ============================================================

\subsection{O programa fork2.c}

O primeiro exemplo da aula coloca \cmd{fork()} dentro de um laço de três iterações. Cada processo filho escreve um valor no vetor \cmd{v}; no final, apenas o processo original soma o vetor e imprime o resultado.

\lstinputlisting[style=arquivo]{\raiz/fork2.c}

Alguns pontos do código antes da análise:

\begin{itemize}
    \item \cmd{calloc} reserva o vetor \textbf{já zerado}. A versão original usava \cmd{malloc}, que não inicializa a memória; a soma lia valores indeterminados.
    \item \cmd{original} guarda o PID do processo que começou o programa. Todos os filhos herdam esse número, mas só um processo tem \cmd{getpid() == original}.
    \item Depois do laço, cada processo espera \textbf{os seus próprios filhos} com \cmd{waitpid(-1, ...)} até receber \cmd{ECHILD}, que significa ``não há mais filhos para coletar''.
\end{itemize}

\subsection{Quantos processos existem?}

O ponto central é que \textbf{os filhos também continuam o laço}. Um filho criado na iteração \cmd{i = 0} executa as iterações \cmd{i = 1} e \cmd{i = 2}, criando os seus próprios filhos.

Vamos acompanhar iteração por iteração, supondo que todas as chamadas tenham sucesso. Chamaremos o processo inicial de \emph{original} e os demais de A a G.

\begin{center}
\begin{tabular}{@{}llc@{}}
\toprule
\textbf{Iteração} & \textbf{Quem chama fork() e quem nasce} & \textbf{Processos no total} \\
\midrule
antes do laço & apenas o original & 1 \\
\cmd{i = 0} & original cria A & 2 \\
\cmd{i = 1} & original cria B; A cria C & 4 \\
\cmd{i = 2} & original cria D; A cria E; B cria F; C cria G & 8 \\
\bottomrule
\end{tabular}
\end{center}

A árvore resultante, com o valor de \cmd{i} em que cada filho nasceu:

\begin{center}
\begin{tikzpicture}[x=1.5cm, y=1.05cm]
    \node[destaque] (p0) at (0,0) {original};
    \node[procf] (a) at (2,0) {A};
    \node[procf] (c) at (3.5,0) {C};
    \node[procf] (g) at (5,0) {G};
    \node[procf] (e) at (3.5,-1) {E};
    \node[procf] (b) at (2,-2) {B};
    \node[procf] (f) at (3.5,-2) {F};
    \node[procf] (d) at (2,-3) {D};
    \draw[seta] (p0) -- node[rotulo, above] {i=0} (a);
    \draw[seta] (a) -- node[rotulo, above] {i=1} (c);
    \draw[seta] (c) -- node[rotulo, above] {i=2} (g);
    \draw[seta] (a.south) |- node[rotulo, pos=0.75, above] {i=2} (e.west);
    \draw[seta] (p0.south) |- node[rotulo, pos=0.8, above] {i=1} (b.west);
    \draw[seta] (b) -- node[rotulo, above] {i=2} (f);
    \draw[seta] (p0.south) ++(-0.3,0) |- node[rotulo, pos=0.8, above] {i=2} (d.west);
\end{tikzpicture}
\end{center}

A cada iteração, o número de processos dobra. Com $n$ iterações:

\begin{itemize}
    \item existem $2^n$ processos ao todo, incluindo o original (aqui, $2^3 = 8$);
    \item $2^n - 1$ deles são novos (aqui, 7);
    \item foram executadas $1 + 2 + 4 + \dots + 2^{n-1} = 2^n - 1$ chamadas de \cmd{fork()}, embora o código tenha \textbf{uma só} chamada escrita.
\end{itemize}

\begin{alerta}[Oito criados não é o mesmo que oito vivos]
A contagem diz quantos processos foram criados ao longo da execução. Não significa que os oito estejam vivos ao mesmo tempo: dependendo do escalonamento, um processo pode terminar antes que outro seja criado.
\end{alerta}

\begin{pergunta}
E se cada filho saísse do laço logo depois de nascer (por exemplo, com um \cmd{break} dentro do \cmd{if (pid == 0)})? Nesse caso só o original completaria as três iterações, criando um filho em cada uma: seriam \textbf{4 processos} ao todo. Em geral, $n + 1$ em vez de $2^n$.
\end{pergunta}

% ============================================================
\section{Cada processo tem a sua memória}\label{sec:memoria}
% ============================================================

\subsection{A previsão}

Antes de executar \cmd{fork2.c}, vale fazer a previsão. Os filhos escrevem \cmd{v[0] = 5}, \cmd{v[1] = 6} e \cmd{v[2] = 7}. Se essas escritas chegassem ao vetor do original, a soma seria $5 + 6 + 7 = 18$. O que o programa imprime?

\begin{lstlisting}[style=terminal]
$ make run-fork
./build/fork2
Soma total no pai: 0
\end{lstlisting}

\subsection{Por que zero}

O original nunca entra no ramo \cmd{if (pid == 0)}: para ele, \cmd{fork()} sempre devolve o PID de um filho. Portanto o original nunca escreve em \cmd{v}. Cada filho escreve na \textbf{sua cópia} do vetor, e essa escrita não alcança a memória de mais ninguém.

\begin{center}
\begin{tikzpicture}
    \node[font=\small\bfseries, anchor=east] at (-0.2,1.6) {Pai antes};
    \foreach \x in {0,1,2} \node[celula] at (\x*0.7+0.35,1.6) {0};
    \draw[seta] (2.3,1.6) -- node[rotulo, above] {fork} (4.2,2.4);
    \draw[seta] (2.3,1.6) -- (4.2,0.8);
    \node[font=\small\bfseries, anchor=west] at (4.3,2.4) {Filho A};
    \node[celula, fill=mLightBrown!25] at (6.1,2.4) {5};
    \node[celula] at (6.8,2.4) {0};
    \node[celula] at (7.5,2.4) {0};
    \node[font=\scriptsize, anchor=west] at (8.0,2.4) {\texttt{v[0] = 0 + 5}};
    \node[font=\small\bfseries, anchor=west] at (4.3,0.8) {Pai};
    \foreach \x in {0,1,2} \node[celula] at (\x*0.7+6.1,0.8) {0};
    \node[font=\scriptsize, anchor=west] at (8.0,0.8) {não entra em \texttt{pid == 0}};
\end{tikzpicture}
\end{center}

Fica mais interessante quando olhamos o vetor de cada um dos oito processos. Um filho herda o vetor \emph{como estava no pai} no momento do \cmd{fork()}, e depois escreve a sua própria posição:

\begin{center}
\begin{tabular}{@{}lllc@{}}
\toprule
\textbf{Processo} & \textbf{Criado por} & \textbf{Vetor ao final} & \textbf{Soma} \\
\midrule
original & (início) & \cmd{[0, 0, 0]} & 0 \\
A & original, \cmd{i = 0} & \cmd{[5, 0, 0]} & 5 \\
B & original, \cmd{i = 1} & \cmd{[0, 6, 0]} & 6 \\
D & original, \cmd{i = 2} & \cmd{[0, 0, 7]} & 7 \\
C & A, \cmd{i = 1} & \cmd{[5, 6, 0]} & 11 \\
E & A, \cmd{i = 2} & \cmd{[5, 0, 7]} & 12 \\
F & B, \cmd{i = 2} & \cmd{[0, 6, 7]} & 13 \\
G & C, \cmd{i = 2} & \cmd{[5, 6, 7]} & 18 \\
\bottomrule
\end{tabular}
\end{center}

A soma 18 existe, mas só na memória de G, que herdou as escritas dos seus ancestrais A e C. Como G não é o original, ele não imprime nada. O original, que imprime, tem o vetor zerado.

\subsection{Mesmo endereço, valores diferentes}

Uma objeção comum: ``mas o endereço de \cmd{v} é o mesmo no pai e no filho''. É verdade, e o programa \cmd{memoria\_independente.c} mostra isso:

\lstinputlisting[style=arquivo]{\raiz/examples/dia05/memoria_independente.c}

\begin{lstlisting}[style=terminal]
$ make run-memoria
./build/memoria-independente
Filho: x = 99 em 0x7ffe72feea88
Pai:   x = 10 em 0x7ffe72feea88
\end{lstlisting}

O endereço impresso é o mesmo, e os valores são diferentes. Não há contradição: o endereço é \textbf{virtual}. O número \cmd{0x7ffe72feea88} é interpretado dentro do espaço de endereçamento de cada processo, e o SO o traduz para posições físicas diferentes no pai e no filho. Mesmo endereço virtual não significa mesma memória.

\subsection{E se o pai esperar os filhos?}

O \cmd{fork2.c} já espera todos os filhos antes de somar, e mesmo assim a soma é zero. Esperar não resolve, porque \cmd{waitpid} não foi feito para transportar dados. Vamos ver o que ele faz.

% ============================================================
\section{Término e coleta: exit e waitpid}
% ============================================================

\subsection{Como um processo termina}

Um processo termina quando retorna de \cmd{main}, quando chama \cmd{exit} ou \cmd{\_exit}, ou quando recebe um sinal que o encerra (por exemplo, \cmd{SIGKILL}). Ao terminar, ele deixa um \textbf{código de saída}: por convenção, \cmd{0} (\cmd{EXIT\_SUCCESS}) indica sucesso e qualquer outro valor indica falha.

\subsection{waitpid}

\begin{lstlisting}[style=codigo]
#include <sys/wait.h>

pid_t waitpid(pid_t pid, int *status, int opcoes);
\end{lstlisting}

\begin{itemize}
    \item \cmd{pid} positivo espera aquele filho específico; \cmd{-1} espera qualquer filho.
    \item Em \cmd{status}, o SO informa como o filho terminou.
    \item Com \cmd{opcoes = 0}, a chamada \textbf{bloqueia} até algum filho terminar.
    \item Devolve o PID do filho coletado, ou \cmd{-1} em caso de erro.
\end{itemize}

O inteiro \cmd{status} é codificado e deve ser lido com macros:

\begin{center}
\begin{tabular}{@{}ll@{}}
\toprule
\textbf{Macro} & \textbf{Pergunta que responde} \\
\midrule
\cmd{WIFEXITED(status)} & O filho terminou normalmente (por \cmd{exit} ou retorno de \cmd{main})? \\
\cmd{WEXITSTATUS(status)} & Se sim, qual foi o código de saída? \\
\cmd{WIFSIGNALED(status)} & O filho foi encerrado por um sinal? \\
\cmd{WTERMSIG(status)} & Se sim, qual sinal? \\
\bottomrule
\end{tabular}
\end{center}

Dois valores de \cmd{errno} aparecem nos exemplos:

\begin{itemize}
    \item \cmd{EINTR}: a espera foi interrompida por um sinal antes de algum filho terminar. Basta chamar \cmd{waitpid} de novo, e é isso que o laço \cmd{do \{ ... \} while} dos exemplos faz.
    \item \cmd{ECHILD}: o processo não tem (mais) filhos para esperar. Em \cmd{fork2.c}, é a condição de saída do laço de coleta.
\end{itemize}

\subsection{Zumbis e órfãos}

Quando um filho termina, o SO não pode descartar tudo: o pai ainda pode querer saber o código de saída. O processo terminado fica então num estado especial, como \textbf{zumbi}, ocupando uma entrada na tabela de processos até o pai coletá-lo com \cmd{wait} ou \cmd{waitpid}. No \cmd{ps}, aparece com estado \cmd{Z}.

\begin{center}
\begin{tikzpicture}[node distance=1.4cm]
    \node[proc] (exec) {em execução};
    \node[destaque, right=of exec] (zumbi) {zumbi};
    \node[proc, right=of zumbi, fill=white] (fim) {removido};
    \draw[seta] (exec) -- node[rotulo, above] {termina} (zumbi);
    \draw[seta] (zumbi) -- node[rotulo, above] {pai coleta} (fim);
\end{tikzpicture}
\end{center}

Se o pai termina antes do filho, o filho fica \textbf{órfão} e é adotado por outro processo, que passa a ser o seu pai e se encarrega de coletá-lo quando ele terminar. Quem adota depende do sistema:

\begin{itemize}
    \item Num servidor ou num terminal sem sessão gráfica, costuma ser o PID 1 (\cmd{init} ou \cmd{systemd}).
    \item Numa sessão gráfica de um Linux atual, costuma ser o \cmd{systemd -{}-user} do usuário. Ele se registra como \textbf{coletor de órfãos} (\emph{subreaper}) e adota os órfãos dos processos da sessão.
\end{itemize}

Por isso, se você imprimir \cmd{getppid()} num órfão no computador do laboratório, provavelmente verá um número diferente de 1. Confira quem é com \cmd{ps -o pid,args -p <número>}.

Um zumbi que dura muito tempo é sinal de um pai que esqueceu de chamar \cmd{waitpid}.

\subsection{Esperar não é comunicar}

\begin{center}
\begin{tabular}{@{}p{0.45\linewidth}p{0.45\linewidth}@{}}
\toprule
\textbf{O que waitpid faz} & \textbf{O que waitpid não faz} \\
\midrule
Espera um filho terminar. & Não copia a memória do filho para o pai. \\
Informa como ele terminou (código ou sinal). & Não cria memória compartilhada. \\
Libera a entrada do zumbi. & Não transporta resultados calculados. \\
\bottomrule
\end{tabular}
\end{center}

\begin{definicao}[Ideia-chave]
\cmd{waitpid} coordena o \textbf{término}. Para levar um resultado do filho ao pai, os bytes precisam ser enviados explicitamente por algum mecanismo do SO.
\end{definicao}

\begin{pergunta}[E se o filho devolvesse o 30 no código de saída?]
Funciona: se o filho chamar \cmd{\_exit(30)}, o pai lê \cmd{WEXITSTATUS(status) == 30}. Mas o código de saída tem só \textbf{8 bits}, ou seja, valores de 0 a 255. Um \cmd{\_exit(300)} chega ao pai como 44 (300 módulo 256). Não cabe nele um número negativo, um \cmd{double}, um texto ou um vetor. Além disso, o valor só chega quando o filho \emph{termina}: um filho que precisa enviar vários resultados ao longo do tempo não tem como fazer isso. O código de saída foi pensado para dizer \emph{como} o processo terminou (0 para sucesso, outro valor para o tipo de falha), e é assim que os nossos exemplos o usam.
\end{pergunta}

% ============================================================
\section{Comunicação entre processos com pipe}\label{sec:pipe}
% ============================================================

\subsection{IPC}

Como a memória de cada processo é isolada, o SO oferece mecanismos para que processos troquem dados. Eles recebem o nome geral de \textbf{IPC} (\emph{Inter-Process Communication}). Pipes, sockets, filas de mensagens e memória compartilhada são exemplos. Nesta aula usamos o mais simples deles: o pipe.

\subsection{O que é um pipe}

\begin{definicao}
Um \textbf{pipe} é um canal de bytes \textbf{unidirecional} mantido pelo núcleo do SO. Ele tem duas pontas: tudo o que é escrito numa ponta pode ser lido, na mesma ordem, na outra.
\end{definicao}

Você já usa pipes no terminal: em \cmd{ps aux | grep sleep}, o shell cria um pipe, liga a saída de \cmd{ps} à ponta de escrita e a entrada de \cmd{grep} à ponta de leitura.

\begin{lstlisting}[style=codigo]
#include <unistd.h>

int pipe(int fd[2]);   /* 0 em sucesso, -1 em erro */
\end{lstlisting}

A chamada preenche o vetor com dois \textbf{descritores de arquivo}:

\begin{itemize}
    \item \cmd{fd[0]}: a ponta de \textbf{leitura};
    \item \cmd{fd[1]}: a ponta de \textbf{escrita}.
\end{itemize}

Uma forma de lembrar: 0 é a entrada padrão (lemos dela) e 1 é a saída padrão (escrevemos nela).

\subsection{Descritores de arquivo}

Um descritor é um número inteiro pequeno que o processo usa para se referir a um recurso aberto: um arquivo, um terminal, um pipe. Cada processo tem a sua \textbf{tabela de descritores}; o número é um índice nessa tabela. Os três primeiros já vêm abertos: \cmd{0} (entrada padrão), \cmd{1} (saída padrão) e \cmd{2} (saída de erros). Por isso \cmd{pipe()} costuma devolver \cmd{3} e \cmd{4}.

Quando o processo chama \cmd{fork()}, o filho recebe uma \textbf{cópia da tabela}. As entradas copiadas apontam para os \emph{mesmos} recursos do núcleo. É isso que permite a comunicação: se o pipe é criado \textbf{antes} do \cmd{fork()}, pai e filho passam a ter referências ao mesmo canal.

\begin{center}
\begin{tikzpicture}[
    tab/.style={draw=mDarkTeal, minimum width=1.5cm, minimum height=0.55cm, font=\ttfamily\footnotesize, fill=white},
    tabx/.style={tab, fill=black!8, text=black!45}]
    % Filho
    \node[font=\small\bfseries] at (0,2.35) {Filho};
    \node[tab] (f0) at (0,1.65) {0 entrada};
    \node[tab, below=0pt of f0] (f1) {1 saída};
    \node[tab, below=0pt of f1] (f2) {2 erros};
    \node[tabx, below=0pt of f2] (f3) {3 fechado};
    \node[tab, below=0pt of f3, fill=mLightBrown!15] (f4) {4 = fd[1]};
    % Pai
    \node[font=\small\bfseries] at (9,2.35) {Pai};
    \node[tab] (p0) at (9,1.65) {0 entrada};
    \node[tab, below=0pt of p0] (p1) {1 saída};
    \node[tab, below=0pt of p1] (p2) {2 erros};
    \node[tab, below=0pt of p2, fill=mLightBrown!15] (p3) {3 = fd[0]};
    \node[tabx, below=0pt of p3] (p4) {4 fechado};
    % Pipe
    \node[draw=mLightBrown, fill=mLightBrown!12, thick, minimum width=3.2cm, minimum height=0.8cm, font=\small]
        (pipe) at (4.5,-0.45) {buffer do pipe};
    \node[font=\scriptsize, color=black!60, below=0.1cm of pipe] {núcleo do SO};
    \draw[setad] (f4.east) -- node[rotulo, above, sloped] {escreve} (pipe.west);
    \draw[setad] (pipe.east) -- node[rotulo, above, sloped] {lê} (p3.west);
\end{tikzpicture}
\end{center}

A figura mostra o estado depois que cada processo fechou a ponta que não usa: o filho fechou a leitura (\cmd{3}) e o pai fechou a escrita (\cmd{4}).

\subsection{A ordem das operações}

\begin{enumerate}
    \item \cmd{pipe(fd)}: o pai cria o canal.
    \item \cmd{fork()}: o filho herda os dois descritores.
    \item Cada processo \textbf{fecha a ponta que não usa}: o filho, que só escreve, fecha \cmd{fd[0]}; o pai, que só lê, fecha \cmd{fd[1]}.
    \item O filho escreve com \cmd{write(fd[1], ...)}; o pai lê com \cmd{read(fd[0], ...)}.
    \item Cada um fecha a ponta que usou quando termina.
    \item O pai coleta o filho com \cmd{waitpid}.
\end{enumerate}

\begin{alerta}[Pipe criado depois do fork não serve]
Se \cmd{pipe()} for chamado depois de \cmd{fork()}, pai e filho criam \emph{dois pipes diferentes}, cada um visível só para quem o criou. Não há canal em comum.
\end{alerta}

\subsection{read e write}

\begin{lstlisting}[style=codigo]
ssize_t write(int fd, const void *buf, size_t n);
ssize_t read(int fd, void *buf, size_t n);
\end{lstlisting}

As duas chamadas trabalham com \textbf{bytes}. \cmd{write} copia até \cmd{n} bytes a partir do endereço \cmd{buf} para o canal; \cmd{read} copia até \cmd{n} bytes do canal para o endereço \cmd{buf}. Ambas devolvem quantos bytes foram de fato transferidos, ou \cmd{-1} em caso de erro.

Para enviar um \cmd{int}, passamos o endereço da variável e o seu tamanho:

\begin{lstlisting}[style=codigo]
write(fd[1], &valor, sizeof valor);   /* no filho */
read(fd[0], &valor, sizeof valor);    /* no pai */
\end{lstlisting}

\begin{definicao}[O ponteiro não é compartilhado]
\cmd{\&valor} no filho e \cmd{\&valor} no pai são endereços de variáveis locais de cada processo. O SO \textbf{copia os bytes} da variável do filho para o buffer do pipe e, depois, do buffer para a variável do pai. As duas variáveis continuam separadas.
\end{definicao}

Enviar os bytes de um \cmd{int} (ou de um \cmd{double}, como em \cmd{pipe2.c}) funciona porque pai e filho rodam na mesma máquina e representam números da mesma forma. Entre computadores diferentes seria preciso combinar um formato, como texto.

\subsection{Bloqueio e fim de arquivo}

O comportamento de \cmd{read} depende do estado do pipe:

\begin{center}
\begin{tabular}{@{}p{0.42\linewidth}p{0.48\linewidth}@{}}
\toprule
\textbf{Situação} & \textbf{O que read faz} \\
\midrule
Há dados no pipe & Devolve os bytes disponíveis, até \cmd{n}. Pode devolver \emph{menos} do que o pedido. \\
Pipe vazio e alguma ponta de escrita ainda aberta & \textbf{Bloqueia}: o processo dorme até chegarem dados. \\
Pipe vazio e todas as pontas de escrita fechadas & Devolve \cmd{0}: \textbf{fim de arquivo} (EOF). \\
\bottomrule
\end{tabular}
\end{center}

Do lado da escrita, também há regras:

\begin{itemize}
    \item O buffer do pipe tem capacidade limitada (64\,KiB por padrão no Linux). Se estiver cheio, \cmd{write} bloqueia até o leitor consumir dados.
    \item Escritas de até \cmd{PIPE\_BUF} bytes (4096 no Linux) são \textbf{atômicas}: não se misturam com escritas de outros processos no mesmo pipe. Um \cmd{int} cabe com folga.
    \item Escrever num pipe sem nenhuma ponta de leitura aberta gera o sinal \cmd{SIGPIPE}, que por padrão encerra o processo. Se o sinal for ignorado, \cmd{write} devolve \cmd{-1} com \cmd{errno = EPIPE}. O \cmd{pipe2.c} ignora \cmd{SIGPIPE} justamente para poder tratar esse erro.
\end{itemize}

\begin{pergunta}[Por que fechar as pontas que não usamos?]
Suponha que o pai esqueça de fechar a sua cópia de \cmd{fd[1]}. Se o filho terminar sem escrever, o pipe fica vazio, mas ainda existe uma ponta de escrita aberta: a do próprio pai. O \cmd{read} do pai então espera para sempre, porque o EOF só chega quando \textbf{todas} as pontas de escrita estão fechadas, inclusive as cópias herdadas.
\end{pergunta}

\subsection{O exemplo pipe\_soma.c}

O filho calcula \cmd{10 + 20} e envia o resultado ao pai:

\lstinputlisting[style=arquivo]{\raiz/examples/dia05/pipe_soma.c}

\begin{lstlisting}[style=terminal]
$ make run-pipe-soma
./build/pipe-soma
Resultado recebido: 30
\end{lstlisting}

\subsubsection{O helper transferir}
Como \cmd{read} e \cmd{write} podem transferir menos bytes do que o pedido, chamá-las uma única vez não garante que o \cmd{int} inteiro foi enviado ou recebido. A função \cmd{transferir} resolve isso com um laço:

\begin{itemize}
    \item \cmd{modo = 1} chama \cmd{write}; \cmd{modo = 0} chama \cmd{read}.
    \item A cada volta, avança o ponteiro \cmd{p} e diminui \cmd{tamanho} pelo número de bytes transferidos, até não restar nada.
    \item Se a chamada foi interrompida por um sinal (\cmd{EINTR}), tenta de novo.
    \item Se \cmd{read} devolve \cmd{0} antes do fim, o escritor fechou o canal cedo demais; isso é tratado como erro.
\end{itemize}

Na atividade prática você vai \emph{usar} esse helper, não reescrevê-lo. O foco é o protocolo: quem fecha qual ponta, quem envia e quem recebe.

\subsubsection{A verificação final do pai}
O pai só imprime o valor se três coisas deram certo: a recepção (\cmd{resultado == 0}), o término normal do filho (\cmd{WIFEXITED}) e o código de saída zero. Aqui \cmd{waitpid} e o pipe trabalham juntos: o pipe trouxe o dado, e \cmd{waitpid} confirmou que o filho terminou bem.

% ============================================================
\section{Vários processos: pipe2.c}
% ============================================================

\subsection{O problema}

O \cmd{pipe2.c} calcula as raízes de uma equação do segundo grau usando seis processos e cinco pipes. O cálculo é pequeno e não fica mais rápido assim; o programa existe para mostrar vários canais funcionando ao mesmo tempo.

O usuário digita \cmd{a}, \cmd{b} e \cmd{delta} (já calculado; o programa não pede \cmd{c}). As raízes são

\[
x_1 = \frac{-b - \sqrt{\Delta}}{2a}, \qquad x_2 = \frac{-b + \sqrt{\Delta}}{2a}.
\]

\subsection{Quem cria quem}

\begin{center}
\begin{tikzpicture}[node distance=0.6cm and 1.2cm]
    \node[destaque] (pai) {Pai original};
    \node[proc, below left=1cm and 1.2cm of pai] (calc) {Calculador\\\scriptsize\texttt{son1}};
    \node[proc, below right=1cm and 1.2cm of pai] (coord) {Coordenador\\\scriptsize\texttt{son5}};
    \node[procf, minimum width=2cm, below=1cm of coord] (f2) {fornece \texttt{2a}};
    \node[procf, minimum width=2cm, left=0.4cm of f2] (f1) {fornece \texttt{-b}};
    \node[procf, minimum width=2cm, right=0.4cm of f2] (f3) {fornece \texttt{delta}};
    \draw[seta] (pai) -- (calc);
    \draw[seta] (pai) -- (coord);
    \draw[seta] (coord) -- (f1);
    \draw[seta] (coord) -- (f2);
    \draw[seta] (coord) -- (f3);
\end{tikzpicture}
\end{center}

\begin{itemize}
    \item O \textbf{pai} lê a entrada, cria os cinco pipes e depois os dois filhos diretos.
    \item O \textbf{calculador} recebe os três componentes, calcula as raízes e as envia ao pai.
    \item O \textbf{coordenador} cria três \textbf{fornecedores}, um para cada componente, e coleta esses filhos.
\end{itemize}

Os valores digitados são lidos pelo pai \emph{antes} dos \cmd{fork()}, então todos os filhos já os herdam na memória. Mesmo assim, os fornecedores enviam os componentes ao calculador pelos pipes: o objetivo é exercitar o protocolo.

\subsection{Por onde os dados passam}

\begin{center}
\begin{tikzpicture}[node distance=0.35cm]
    \node[procf, minimum width=2.1cm] (f1) {\texttt{-b}};
    \node[procf, minimum width=2.1cm, below=of f1] (f2) {\texttt{2a}};
    \node[procf, minimum width=2.1cm, below=of f2] (f3) {\texttt{delta}};
    \node[proc, right=3.2cm of f2, minimum height=1.6cm] (calc) {Calculador};
    \node[destaque, right=3.2cm of calc] (pai) {Pai};
    \draw[setad] (f1.east) -- node[rotulo, above, sloped] {\texttt{canais[0]}} (calc.west);
    \draw[setad] (f2.east) -- node[rotulo, above] {\texttt{canais[1]}} (calc.west);
    \draw[setad] (f3.east) -- node[rotulo, below, sloped] {\texttt{canais[2]}} (calc.west);
    \draw[setad] ([yshift=0.25cm]calc.east) -- node[rotulo, above] {\texttt{canais[3]}: x1} ([yshift=0.25cm]pai.west);
    \draw[setad] ([yshift=-0.25cm]calc.east) -- node[rotulo, below] {\texttt{canais[4]}: x2} ([yshift=-0.25cm]pai.west);
\end{tikzpicture}
\end{center}

\begin{alerta}[Árvore de criação e fluxo de dados são coisas diferentes]
Os fornecedores são filhos do coordenador, mas não enviam nada a ele: enviam ao calculador, que é seu ``tio''. Quem pode usar um pipe é quem tem um descritor dele, e todos herdaram os cinco canais do pai original. A árvore diz quem criou quem; os pipes dizem quem conversa com quem.
\end{alerta}

\subsection{Cada processo fica só com as pontas que usa}

A função \cmd{fechar\_exceto} fecha todos os descritores dos cinco pipes, exceto os indicados. A tabela resume o que cada processo mantém aberto:

\begin{center}
\begin{tabular}{@{}lll@{}}
\toprule
\textbf{Processo} & \textbf{Lê de} & \textbf{Escreve em} \\
\midrule
Pai & \cmd{canais[3]}, \cmd{canais[4]} & nenhum \\
Calculador & \cmd{canais[0]}, \cmd{canais[1]}, \cmd{canais[2]} & \cmd{canais[3]}, \cmd{canais[4]} \\
Coordenador & nenhum & \cmd{canais[0..2]}, só até criar os fornecedores \\
Fornecedor \cmd{i} & nenhum & \cmd{canais[i]} \\
\bottomrule
\end{tabular}
\end{center}

O coordenador mantém as três pontas de escrita apenas para que os fornecedores as herdem; logo depois de criá-los, fecha todas. Se não fechasse, o calculador nunca veria EOF nesses canais caso um fornecedor falhasse.

\subsection{Os trechos principais}

O código completo está em \cmd{pipe2.c}. Os trechos abaixo concentram as ideias da aula.

\lstinputlisting[style=arquivo, firstline=99, lastline=120, firstnumber=99]{\raiz/pipe2.c}

O calculador fecha tudo o que não usa, recebe os três valores, calcula e envia as raízes. Os fornecedores são criados num laço dentro do coordenador:

\lstinputlisting[style=arquivo, firstline=122, lastline=152, firstnumber=122]{\raiz/pipe2.c}

Note o \cmd{for} interno no filho: cada fornecedor fecha as pontas de escrita dos \emph{outros} dois canais e fica só com a sua.

\subsection{Execução}

\begin{lstlisting}[style=terminal]
$ make run-pipe
./build/pipe2
>>> Digite tres numeros separados por espaco e tecle Enter: a b delta
>>> Exemplo: 1 -3 1   (delta ja calculado; o programa nao pede c)
a b delta: 1 -3 1
Raizes: x1 = 1; x2 = 2
\end{lstlisting}

Conferindo: $-b = 3$, $2a = 2$ e $\sqrt{\Delta} = 1$. Então $x_1 = (3 - 1)/2 = 1$ e $x_2 = (3 + 1)/2 = 2$. Com a entrada \cmd{1 -2 0}, o programa mostra uma raiz dupla ($x_1 = x_2 = 1$).

% ============================================================
\section{Atividade prática}
% ============================================================

\subsection{O que fazer}

Abra \cmd{examples/dia05/pipe\_soma\_atividade.c}. Ele é o \cmd{pipe\_soma.c} com seis lacunas, marcadas em laranja abaixo:

\begin{lstlisting}[style=codigo]
if (pid == 0) {
    close(@/* 1: ponta que o filho nao usa */@);
    valor = @/* 2: calculo */@;
    int resultado = @/* 3: enviar valor pelo helper transferir */@;
    close(fd[1]);
    _exit(resultado == 0 ? EXIT_SUCCESS : EXIT_FAILURE);
}
close(@/* 4: ponta que o pai nao usa */@);
int resultado = @/* 5: receber valor pelo helper transferir */@;
close(@/* 6: ponta usada pelo pai */@);
\end{lstlisting}

Complete as lacunas para que:

\begin{itemize}
    \item o filho feche a sua ponta de leitura, calcule \cmd{10 + 20}, envie o inteiro e feche a ponta de escrita;
    \item o pai feche a sua ponta de escrita, receba o inteiro, feche a ponta de leitura e imprima o valor.
\end{itemize}

Lembre a assinatura do helper: \cmd{transferir(descritor, endereço do dado, quantidade de bytes, modo)}, com modo \cmd{1} para enviar e \cmd{0} para receber.

\subsection{Compilar e testar}

\begin{lstlisting}[style=terminal]
$ gcc -std=c11 -Wall -Wextra -Wpedantic -Werror \
      examples/dia05/pipe_soma_atividade.c -o /tmp/pipe-soma-aluno
$ /tmp/pipe-soma-aluno
Resultado recebido: 30
\end{lstlisting}

Antes de preencher, o esqueleto não compila. Isso é proposital.

\subsection{Lista de conferência}

\begin{itemize}
    \item Cada processo fecha \textbf{exatamente uma} ponta antes de usar o canal, e é a ponta que ele não usa.
    \item O filho usa a ponta de escrita; o pai usa a de leitura.
    \item O tamanho passado ao helper é o tamanho do dado (\cmd{sizeof valor}), não o tamanho de um ponteiro.
    \item O endereço passado é o da variável (\cmd{\&valor}), não o seu valor.
    \item Você consegue explicar por que escrever em \cmd{valor} no filho, sem o pipe, não mudaria \cmd{valor} no pai.
\end{itemize}

\subsection{Perguntas para responder em dupla}

\begin{enumerate}
    \item Por que o pipe precisa ser criado antes do \cmd{fork()}?
    \item Se o pai executar \cmd{read} antes de o filho escrever, isso é um erro? Quando o pai espera? Quando encontra EOF?
    \item Por que os dois processos fecham pontas diferentes?
\end{enumerate}

\subsection{Erros comuns}

\begin{center}
\begin{tabular}{@{}p{0.36\linewidth}p{0.56\linewidth}@{}}
\toprule
\textbf{Se você pensou\ldots} & \textbf{Reveja\ldots} \\
\midrule
``O pai soma 18.'' & Em qual processo cada atribuição acontece? Desenhe as memórias separadas. \\
``São quatro processos.'' & Siga um filho na próxima iteração: ele não sai do laço. \\
``waitpid compartilha memória.'' & Separe esperar o término de transportar os bytes. \\
``read sempre recebe tudo.'' & O retorno informa quantos bytes chegaram. \\
``Pipe vazio sempre dá EOF.'' & Alguma ponta de escrita ainda está aberta? \\
``Mesmo endereço, mesma memória.'' & O endereço é virtual e interpretado em cada processo. \\
\bottomrule
\end{tabular}
\end{center}

% ============================================================
\section{Resumo}
% ============================================================

\begin{center}
\begin{tabular}{@{}lp{0.68\linewidth}@{}}
\toprule
\textbf{Chamada} & \textbf{Papel} \\
\midrule
\cmd{fork()} & Cria um processo filho com uma cópia independente da memória e dos descritores do pai. Devolve 0 ao filho, o PID do filho ao pai e -1 em caso de falha. \\
\cmd{getpid()} / \cmd{getppid()} & PID do próprio processo / PID do pai. \\
\cmd{\_exit(codigo)} & Encerra o processo imediatamente, sem esvaziar buffers de \cmd{stdio}. \\
\cmd{waitpid(pid, \&st, 0)} & Espera um filho terminar e coleta o seu estado. Coordena o término; não transporta dados. \\
\cmd{pipe(fd)} & Cria um canal unidirecional: \cmd{fd[0]} lê, \cmd{fd[1]} escreve. \\
\cmd{write} / \cmd{read} & Copiam bytes para o canal e do canal. Podem transferir menos que o pedido; \cmd{read} bloqueia com o canal vazio e devolve 0 no EOF. \\
\cmd{close(fd)} & Fecha um descritor. Fechar todas as pontas de escrita é o que gera EOF para o leitor. \\
\bottomrule
\end{tabular}
\end{center}

\begin{definicao}[Em três frases]
\cmd{fork()} cria processos com memórias independentes. \cmd{waitpid} coordena o término e coleta os filhos. O pipe transporta os dados entre eles.
\end{definicao}

\begin{exemplo}[Próxima aula: 7/10]
Threads de um mesmo processo \textbf{compartilham} a memória. Aí o problema se inverte: os dados chegam sem esforço, mas é preciso controlar quem acessa o quê e quando. Esse é o assunto de região crítica, Peterson e semáforos.
\end{exemplo}

% ============================================================
\section{Exercícios de fixação}
% ============================================================

\begin{enumerate}
    \item Em \cmd{fork2.c}, troque \cmd{tamanho} para 4. Quantos processos existem ao todo? Quantas chamadas de \cmd{fork()} são executadas?
    \item Quantas vezes o programa abaixo imprime \cmd{oi}? Desenhe a árvore.
\begin{lstlisting}[style=codigo]
int main(void) { fork(); fork(); printf("oi\n"); return 0; }
\end{lstlisting}
    \item O programa abaixo imprime \cmd{AA} no terminal, embora haja um único \cmd{printf}. Explique.
\begin{lstlisting}[style=codigo]
int main(void) { printf("A"); fork(); return 0; }
\end{lstlisting}
    \item Na tabela da Seção~\ref{sec:memoria}, qual é o vetor do processo E e por quê?
    \item Em \cmd{pipe\_soma.c}, o que aconteceria se o pai não fechasse \cmd{fd[1]} e o filho terminasse com falha antes de escrever?
    \item Para que pai e filho conversem nos dois sentidos (pergunta e resposta), quantos pipes são necessários? Por quê?
    \item Por que o filho chama \cmd{\_exit} em vez de \cmd{exit} ou \cmd{return}?
    \item Por que \cmd{pipe2.c} ignora o sinal \cmd{SIGPIPE}?
\end{enumerate}

\subsection*{Respostas}

\begin{enumerate}
    \item $2^4 = 16$ processos ao todo, 15 novos; também 15 chamadas executadas ($1 + 2 + 4 + 8$).
    \item Quatro vezes. O primeiro \cmd{fork()} gera 2 processos; cada um executa o segundo \cmd{fork()}, gerando 4. Os quatro chegam ao \cmd{printf}.
    \item O \cmd{A} fica no buffer de \cmd{stdout}, pois não há quebra de linha. O \cmd{fork()} copia o buffer. Cada processo, ao retornar de \cmd{main}, esvazia a sua cópia, e o \cmd{A} aparece duas vezes. Um \cmd{fflush(stdout)} antes do \cmd{fork()} evita a duplicação.
    \item \cmd{[5, 0, 7]}. E foi criado por A na iteração \cmd{i = 2}. Herdou o vetor de A, que já tinha \cmd{v[0] = 5}, e escreveu \cmd{v[2] = 7}.
    \item O \cmd{read} do pai bloquearia para sempre. O pipe estaria vazio, mas a ponta de escrita do próprio pai continuaria aberta, e o EOF nunca chegaria. O pai também nunca chegaria ao \cmd{waitpid}.
    \item Dois. O pipe é unidirecional: um canal leva dados do pai ao filho e o outro, do filho ao pai. Usar um só canal nos dois sentidos faria um processo poder ler o que ele mesmo escreveu.
    \item \cmd{exit} e \cmd{return} esvaziam os buffers de \cmd{stdio} herdados do pai (podendo duplicar saídas) e executam funções registradas pelo pai. \cmd{\_exit} encerra só o filho, sem esses efeitos.
    \item Se o leitor de um canal terminar antes, uma escrita nesse canal geraria \cmd{SIGPIPE} e encerraria o processo sem chance de tratamento. Ignorando o sinal, \cmd{write} devolve \cmd{-1} com \cmd{EPIPE}, e o programa trata a falha normalmente.
\end{enumerate}

% ============================================================
\section*{Para saber mais}
\addcontentsline{toc}{section}{Para saber mais}
% ============================================================

\begin{itemize}
    \item Páginas de manual do Linux: \cmd{man 2 fork}, \cmd{man 2 waitpid}, \cmd{man 2 pipe} e \cmd{man 7 pipe}. Também disponíveis em \url{https://man7.org/linux/man-pages/}.
    \item TANENBAUM, A. S.; BOS, H. \emph{Sistemas Operacionais Modernos}. 4. ed. São Paulo: Pearson, 2016. Capítulo 2: processos e threads.
    \item SILBERSCHATZ, A.; GALVIN, P. B.; GAGNE, G. \emph{Fundamentos de Sistemas Operacionais}. 9. ed. Rio de Janeiro: LTC, 2015. Capítulo 3: processos, incluindo comunicação entre processos e pipes.
\end{itemize}

Em outras edições desses livros, o assunto continua nos capítulos iniciais sobre processos, mas a numeração pode mudar.

\begin{itemize}
    \item KERRISK, M. \emph{The Linux Programming Interface}. San Francisco: No Starch Press, 2010. Capítulos 24 a 26 (criação, término e espera de processos) e 44 (pipes). Referência detalhada, em inglês, escrita pelo mantenedor das páginas de manual do Linux.
\end{itemize}

\end{document}
